\chapter{Juarrero and Deacon: Framework Comparisons}
\label{app:juarrero_deacon}

\noindent This appendix situates the process ontology of this monograph relative
to two major contemporary accounts of emergence, constraint causation, and
intentionality: Alicia Juarrero's constraint-based dynamics and Terrence Deacon's
absential organization. Both influence the framework developed here; neither
is adopted wholesale. The differences are philosophical load-bearing rather
than terminological.

\section{Juarrero: Context-Sensitive Constraints}

Juarrero's central contribution in \emph{Dynamics in Action} \citep{juarrero1999} is the
distinction between first-order and second-order constraints. First-order constraints
are boundary conditions that limit the range of possible system states directly.
Second-order constraints are context-sensitive: they alter the probability
distribution over future states rather than eliminating possibilities outright.
Juarrero argues that intentional behavior is best understood as second-order
constraint causation: the meaning of an action is determined by the context that
constrains which actions are admissible, not by the intrinsic properties of the
action itself.

\textbf{Points of convergence with the present framework:}

The trajectory-admissibility framework is continuous with Juarrero's account.
Admissible trajectories are exactly the states that survive the imposition of
both first-order and second-order constraints. The historical accumulation of
constraints---the decreasing filtration $\mathcal{A}(\mathcal{C}_0) \supseteq
\mathcal{A}(\mathcal{C}_1) \supseteq \ldots$---is a formalization of what
Juarrero means by context-sensitive constraint accumulation. The \RSVP entropy
field $S$ functions as a measure of residual second-order constraint: regions
of high $S$ are regions where the contextual constraints have not yet settled
into a stable pattern.

\textbf{Points of divergence:}

Juarrero's account sometimes approaches a strong reading on which second-order
constraints are themselves causally productive---on which the context does not
merely filter but actively contributes causal force to the outcomes it shapes.
The present framework takes a more deflationary position: constraints are
eliminative rather than productive. They remove inadmissible trajectories but
do not add new causal forces. Apparent goal-directedness is a consequence of
what has been eliminated, not of any positive force directed toward a target.
This is a subtle but important difference in the direction of causation.

\section{Deacon: Absential Organization}

Deacon's account in \emph{Incomplete Nature} \citep{deacon2011} introduces the concept of
\emph{absence} as a causally efficacious feature of complex systems. The central
claim is that teleodynamic organization---the kind that characterizes living
systems and minds---is constituted by the causal influence of what is not present:
potential states toward which the system is organized, constraints that do not
currently exist but that the system's organization tends to produce.

Deacon develops a hierarchy of emergent dynamics:
\begin{itemize}
  \item \textbf{Homeodynamics}: simple attractor dynamics in physical systems.
  \item \textbf{Morphodynamics}: self-organizing systems that develop global
    regularities from local interactions.
  \item \textbf{Teleodynamics}: systems whose organization is constituted by
    reference to absent or potential states.
\end{itemize}

\textbf{Points of convergence with the present framework:}

The teleodynamics level is genuinely important for understanding the kind of
coherence exhibited by living systems and minds. Deacon is right that adequate
accounts of intentionality require more than simple homeodynamics or
morphodynamics. The \KES map's successor hypothesis space $H_{t+1}$ is exactly
an absent state that constrains present dynamics: the system's current trajectory
is partly constituted by which future hypothesis spaces remain admissible.

\textbf{Points of divergence:}

The present framework declines to grant absence ontological primacy. On Deacon's
strongest reading, not-yet-realized states exert genuine causal influence on
current dynamics---there is a kind of backward causation or final causation
operating through the organization of living systems. The present framework
explains the same phenomena without this commitment. The \KES successor
hypothesis space $H_{t+1}$ constrains current dynamics not because it is
causally efficacious as an absence but because the current constraint set
$\mathcal{C}_t$ has already eliminated the trajectories that would make
$H_{t+1}$ unreachable. The direction of causation is entirely forward: prior
constraints eliminate future possibilities; the future does not pull the present.

Formally: the filtration $\mathcal{A}(\mathcal{C}_0) \supseteq
\mathcal{A}(\mathcal{C}_1) \supseteq \ldots$ produces apparent goal-directedness
through elimination alone. No backward causation is required. The system appears
to be directed toward the configurations that survive the filtration because all
alternatives have been progressively eliminated. This is sufficient to generate
the phenomena Deacon attributes to teleodynamics without requiring absence to be
causally efficacious.

\section{The Constraint-Final Causation Distinction}

The formal distinction between constraint causation and final causation is:

\begin{definition}{Constraint Causation}
A system exhibits \emph{constraint causation} if its directional structure
is fully explained by the accumulation of prior constraints that eliminate
inadmissible trajectories. The causal influence flows entirely forward:
past constraint events eliminate future possibilities, producing a
monotonically narrowing admissibility manifold.
\end{definition}

\begin{definition}{Final Causation}
A system exhibits \emph{final causation} if a not-yet-realized future state
$s^*$ exerts causal influence on current system dynamics, such that the
dynamics would be different if $s^*$ were different, holding all prior
history constant.
\end{definition}

The present framework adopts constraint causation throughout and declines final
causation. \claimC This position is consistent with thermodynamics, bounded
computation, and path dependence. It produces all the observed phenomena
(apparent purposiveness, anticipatory behavior, structural persistence) without
positing causation that runs backward through time or that requires future
states to exist ontologically before they are realized.

\section{Where Deacon's Framework Remains Valuable}

Despite the divergence on final causation, Deacon's hierarchical emergence
framework provides useful vocabulary for distinguishing levels of organization.
The homeodynamic/morphodynamic/teleodynamic hierarchy maps onto the trajectory
framework as follows:

Homeodynamic systems navigate a fixed admissibility manifold whose topology
is specified in advance. Morphodynamic systems modify their admissibility manifold
through self-organization, but the modification rules are themselves fixed.
Teleodynamic systems---and this is the genuinely novel level---modify both their
admissibility manifold and their modification rules through ongoing interaction
with the outcomes of their prior modifications.

The Yarncrawler framework captures exactly this last level: the self-rewrite
endofunctor $\mathsf{R}: \mathcal{M} \to \mathcal{M}$ modifies the manifold,
while the RC-CLIO repair cycle modifies the repair rules themselves. This is
teleodynamic organization in Deacon's sense, achieved through forward-only
constraint causation rather than through final causation.
